\documentclass[runningheads,a4paper]{llncs}
\usepackage[T1]{fontenc}
\usepackage{graphicx}
\usepackage{booktabs}
\usepackage{amsmath}
\usepackage{url}
\usepackage{xurl}
\setlength{\emergencystretch}{3em}
\providecommand{\doi}[1]{\url{https://doi.org/#1}}

% ---- Anonymous-review switch: \reviewtrue hides author, affiliation, email ----
\newif\ifreview
\reviewfalse

\begin{document}

\title{Compiler-Guided Proof Repair in Lean 4: A Controlled Local-Model Pilot on Fixed Theorem Types}
\titlerunning{Compiler-Guided Lean 4 Proof Repair Pilot}
\ifreview
  \author{Anonymous author}
  \authorrunning{Anonymous author}
  \institute{Anonymous institution}
\else
  \author{Sadia Peerzada}
  \authorrunning{S. Peerzada}
  \institute{Aligarh Muslim University, Aligarh, India\\
  \email{sadiapeerzada2006@gmail.com}}
\fi
\maketitle

\begin{abstract}
Does showing a compiler diagnostic to a small local model produce more verified
Lean~4 proofs than simply allowing more attempts? We compare one-shot generation (A),
blind repair (B) and compiler-guided repair (C) on fixed, AI-reviewed theorem types,
using shared initial candidates, at most three repairs per branch and one
hash-pinned 4-bit Qwen2.5-7B-Instruct model on 60 held-out targets. Strict
verification of the original type succeeds for A on 0/60, B on 1/60 and C on 1/58
observed targets (two C trajectories are unavailable and are not counted as failures).
On the 58 complete pairs both repair arms verify the same single theorem, so no
success advantage is observed; zero discordance cannot show equivalence, and a
paired-score interval of $[-6.83,+6.83]$ points remains compatible with the data.
On the 58 complete B/C pairs, guided repair incurred approximately 2.2 times the mean per-theorem input-token consumption of blind repair (4,353.5 versus 1,980.9 tokens), without additional verified success.
In an exploratory post-hoc analysis, byte-identical retries were less frequent under C
than B (70/172 = 40.7\% vs.\ 90/178 = 50.6\%), yet these changed strings yielded no
additional verified theorem, and a 15-problem development placebo is ambiguous.
The contribution is an auditable application of a controlled blind-versus-guided repair comparison, not a new method; it shows that string change, error change and kernel verification are distinct measurements. One model,
one seed, a single solved target and AI-assisted reference review limit
generalization; autonomous formalization is not evaluated.
\keywords{Lean 4 \and proof repair \and compiler diagnostics \and local inference \and negative results}
\end{abstract}

%==================================================================
\section{Introduction and Research Questions}
A theorem prover can reject a generated proof for many distinct reasons:
malformed syntax, unavailable identifiers, ill-typed expressions, unmet proof
goals, or an incorrect mathematical argument. Compiler diagnostics expose some of
these obstacles, which makes feedback-conditioned rewriting an appealing
inference-time intervention. However, extra model calls can change a candidate
even when the model never sees the error. Distinguishing feedback from additional attempts calls for a blind-retry comparator with the same repair opportunities, rather than comparison against a single generation. Such a comparator controls for extra repair opportunities and shared initial candidates, but it does not isolate diagnostic content from all other condition-specific differences, including prompt length and random sampling.

A second difficulty concerns what counts as improvement. A repair loop can report a
changed proof string, a vanished error message or a different error category, none
of which implies that the original theorem is now proved. We therefore treat
compiler-reported change and strict verification as separate measurements and make
only the latter primary.

We study a deliberately narrow question. With the theorem type and the initial
failed candidate fixed, does real compiler feedback help a small local model
produce more kernel-accepted proofs than blind repair? The theorem trajectory,
not each repair, is the unit of analysis. We ask:
\begin{itemize}
\item \textbf{RQ1 (primary, frozen).} Does compiler-guided repair (C) verify more
fixed targets than equally budgeted blind repair (B), given the same initial candidate?
\item \textbf{RQ2 (secondary, frozen).} How much do B and C add over one-shot
generation (A), and for which initial error categories?
\item \textbf{RQ3 (exploratory, post hoc).} Do diagnostics change \emph{what the
model does} (e.g.\ how often it repeats its previous candidate), and does that
change depend on the diagnostic's content?
\end{itemize}
Here ``equally budgeted'' in RQ1 refers to the cap of three repair attempts, not to equal input tokens or computation. Even an inconclusive primary result is informative for evaluation design: when
success is this rare, a process metric such as repetition can move while verified
success does not, and a study that reported only the process metric would have
suggested a benefit that the endpoint does not support. Three distinctions organize
the paper. A changed proof string is not a better proof. Absence of an observed
advantage is not equivalence, especially when almost every trajectory fails. And
success on a fixed, reviewed Lean type cannot validate an autonomous English-to-Lean
pipeline, since the model never chooses the theorem type.

\smallskip\noindent\textbf{Contributions.} This paper proposes no new prover, algorithm or benchmark, and does not claim to have introduced feedback-based repair, blind retries or paired experimental controls. It offers (1)~an auditable application of a controlled blind-versus-guided repair comparison in a resource-constrained local-model setting, with shared initial candidates, bounded repairs, separate random streams and strict verification;
(2)~transparent primary negative or inconclusive evidence, with unavailable
trajectories and diagnostic-prompt cost reported rather than imputed; and
(3)~an explicitly exploratory, traceable characterization of candidate behavior and
failure modes (frozen error taxonomy, repetition analysis, mismatched-feedback
placebo, authentic cases). We do not claim that compiler feedback is ineffective in general.

%==================================================================
\section{Related Work}
\paragraph{Lean and benchmarks.} Lean~4 provides kernel-checked proof
infrastructure~\cite{demoura2021lean4} and Mathlib the formal library used
here~\cite{mathlib2020}. miniF2F~\cite{zheng2022minif2f} targets formal
olympiad-level problems, ProofNet~\cite{azerbayev2023proofnet} relates informal
and formal undergraduate mathematics, and LeanDojo~\cite{yang2023leandojo}
studies retrieval-augmented proving. Draft, Sketch, and
Prove~\cite{jiang2023dsp} and DeepSeek-Prover~\cite{xin2024deepseekprover} use
informal sketches or large synthetic training data. These systems differ in task
construction, search, premise access and training, so their success rates are not
comparable with ours.

\paragraph{Feedback versus blind sampling.} Iterative refinement with execution or
self-generated feedback is studied for code
(Self-Debug~\cite{chen2024selfdebug}), general generation
(Self-Refine~\cite{madaan2023selfrefine}) and agents
(Reflexion~\cite{shinn2023reflexion}). Huang et al.~\cite{huang2024selfcorrect}
report that models often fail to correct reasoning without reliable external
feedback, which motivates asking whether an \emph{external verifier signal},
rather than extra sampling, drives gains. The Lean compiler supplies such a signal;
our blind-retry arm B controls for extra attempts, although it does not isolate diagnostic content from other condition-specific differences.

\paragraph{Trained proof-repair systems.} Baldur~\cite{first2023baldur} performs
whole-proof generation and repair for Isabelle with fine-tuned models, using error
messages in repair. APRIL~\cite{wang2026april} trains Lean models on about 260,000
supervised repair examples to repair proofs from compiler feedback. Our setting
differs in kind: inference-time feedback to one \emph{unfine-tuned} local instruction
model~\cite{qwen2024report}, with a paired blind-retry control and a separate
post-hoc mismatched-feedback control. We do not compare raw success rates with
these systems, because proof assistant, task, model training and budget all
differ. The contribution is an auditable application of a controlled blind-versus-guided repair comparison in a resource-constrained local-model setting, with explicit accounting of verified outcomes, incomplete trajectories, exploratory candidate behavior and diagnostic-prompt costs. We do not claim that this comparator is the first of its kind: no systematic prior-art review was performed, and the supplied references already establish substantial work on feedback-conditioned proof repair. We test neither trained
repair models, retrieval, search trees, nor stronger foundation models.

\paragraph{Verifier-feedback component controls.} Li and Li~\cite{li2026hindsight}
study Lean proof repair with controlled failed-predecessor exposure,
correct versus mismatched diagnostics, instruction-only retrying, and
token-budget comparisons. Their Qwen and specialized-prover experiments
find that feedback can affect behavior without establishing an advantage
over independent resampling in the evaluated settings. Our pilot does not
replicate their interventions or match their statistical power: it studies
bounded multi-step trajectories from shared initial candidates on fixed
theorem types using one 4-bit local instruction model, while separately
accounting for missing guided outcomes, verified success, surface edits
and resource usage. We therefore claim neither priority for blind controls
nor a comparable effect estimate.

\paragraph{Semantic-faithfulness limits.} Work on faithfulness beyond
compilation~\cite{zhang2026beyond} separates proof synthesis for an existing type
from the correctness of an automatically formalized statement. Our outcome
definition keeps that boundary.

%==================================================================
\section{Study Design}
\paragraph{Pipeline and conditions.} The design matches repair
\emph{opportunities}, not realized tokens: C's diagnostics lengthen prompts, and B
was executed before C without randomization. The blind-retry comparator controls for additional repair opportunities and shared initial candidates, but the comparison does not isolate diagnostic content from all other condition-specific differences, including prompt length and random sampling. For each frozen English/type pair the
pipeline renders an initial prompt, produces one candidate, extracts the Lean
block and verifies it against the reference type. A, B and C share the initial
candidate's exact bytes and SHA-256. Repairs start independent conversations: B
receives the theorem, its current candidate and fixed blind-feedback text; C
receives the same template plus structured diagnostics from the immediately
preceding failed compilation. Neither branch sees the other's candidates. B and C
each allow up to three repairs (at most seven logical generation intents per
problem including the shared initial call); branches stop on verified success.
C is never given fabricated diagnostics, and unavailable diagnostics halt that repair.

\begin{figure}[tbp]
\centering
\includegraphics[width=\linewidth,height=.29\textheight,keepaspectratio]{fig1.png}
\caption{Schematic overview of fixed-theorem proof generation and bounded repair.
The study supplies pre-reviewed Lean theorem types, rather than evaluating
autonomous English-to-Lean translation. The three conditions share an initial
candidate: A uses no repair, B permits three blind retries, and C permits
three retries informed by genuine compiler diagnostics.}
\label{fig:study-pipeline}
\end{figure}

\paragraph{Verification.} A string counts as a proof only if the strict checker
accepts the exact elaborated target type and its axiom dependencies. It rejects
placeholders and unauthorized declarations, permits only \texttt{propext},
\texttt{Classical.choice} and \texttt{Quot.sound}, and rejects \texttt{sorryAx}.
A reference placeholder is used only for type comparison in a separate namespace.
Source-token and declaration guards run before compilation; guard or axiom-policy
rejections are not relabelled as compiler errors. Context-limit failures and
interrupted generations remain explicit missing outcomes. The outcome is
kernel-oriented acceptance for the fixed type, not semantic accuracy of the English.

\paragraph{Model and decoding.} The model is
\path{mlx-community/Qwen2.5-7B-Instruct-4bit}, run offline with
MLX~\cite{qwenmlx,mlxlm}; snapshot
\path{c26a38f6a37d0a51b4e9a1eb3026530fa35d9fed}; 4-bit, group size 64; MLX
0.32.4, MLX-LM 0.32.2, Transformers 5.21.0. Decoding: temperature 0.2, top-$p$ 1,
$\le$2,000 output tokens, 8,192-token context cap, one primary seed (2027), on an
Apple M4 with 16\,GiB. Per-call seeds are the first unsigned 32-bit word of
SHA-256 over base seed, phase/problem/condition/step context and exact prompt
hash, so B's stopping behaviour does not shift C's random stream. Weights are
released during the Lean check and reloaded from the same verified bytes,
preventing KV-state carryover at the cost of loading time. The environment is
Lean 4.32.0 with Mathlib commit \path{81a5d257c8e410db227a6665ed08f64fea08e997};
there is no hosted-inference fallback.

\paragraph{Freeze and execution.} Development pilots exposed Lean~3-style syntax
and high memory use; before any held-out call, the protocol added the same Lean~4
syntax instruction to all arms and fixed the memory schedule (implementation
development, not test tuning). A separate three-case pre-test pilot passed the
provenance/control audit and a resume invocation generated no new calls; it must
not be confused with the later 15-problem post-hoc placebo. No model, prompt,
dataset, seed or budget changed after held-out execution began. B ran before C
(not randomized). Each completed call has a durable intent, raw response and token
stream with verification checkpoints; resumption reuses stored responses and
incomplete intents consume their logical call.

%==================================================================
\section{Dataset and Validity}
The internally frozen, provenance-screened dataset has 75 English/theorem-type
pairs: 15 development and 60 held-out test statements in five and 54 curated
families respectively, with no family overlap (composition, exclusions and
licensing in Appendix~A). The split and exact bytes were frozen with
cryptographic hashes before held-out generation, but not publicly preregistered.

Four properties qualify what the dataset can support. First, reference types were
selected using AI-assisted assessment; independent human labels number zero. The 75
retained references were labelled faithful under the adopted AI eligibility policy,
which is not an independently validated measure of semantic accuracy. Second, 52 of
the 60 test statements and 7 of the 15 development statements are adaptations of
formal-library statements, so results are not contamination-free. Third, difficulty
grades are curated labels, not objectively validated, and 57 of 60 test targets
carry grade 2, so the test set has little recorded difficulty variation. Fourth,
curated families are human groupings, not independence guarantees. The study
therefore evaluates model attempts \emph{conditional on} the fixed types; it does
not estimate semantic-review accuracy or end-to-end autoformalization.

%==================================================================
\section{Results and Statistical Analysis}
\paragraph{One solved theorem.} Across all 60 held-out targets exactly one distinct
theorem, \texttt{lib\_a\_exp\_add}, is verified, and it is solved by both B and C.
One-shot generation (A) solves none. Across the 58 complete paired trajectories (52 curated families), B and C verify the same single theorem, while the remaining 57 targets fail under both conditions; no target is verified only under B and none only under C. With zero discordant verified-success outcomes, the observed paired difference is zero; however, the sparse data cannot establish equivalence or precisely determine the population effect.

\begin{table}[tbp]
\caption{Observed verification outcomes and known completed-response resources.
C has two unavailable outcomes (primary rate 1/58). Repair means use the 60-case
denominator. Resource totals include the shared initial in each counterfactual row
and impute nothing for missing calls.}
\label{tab:main}
\centering\footnotesize\setlength{\tabcolsep}{3pt}
\resizebox{\linewidth}{!}{\begin{tabular}{lcccccccc}
\toprule
Cond. & Ver./obs. & Miss. & Rate & Wilson 95\% & Repairs (mean) & Tokens in/out & Call s & Lean s\\
\midrule
A & 0/60 & 0 & 0.00\% & [0.00, 6.02] & 0 (0.00)    & 12,929/27,862   & 1908.6 & 2361.9\\
B & 1/60 & 0 & 1.67\% & [0.29, 8.86] & 178 (2.97)  & 122,029/83,485  & 8120.5 & 8923.9\\
C & 1/58 & 2 & 1.72\% & [0.31, 9.14] & 172 (2.87)  & 252,935/74,778  & 9870.8 & 8906.3\\
\bottomrule
\end{tabular}}
\end{table}

A verifies 0/60 (0.00\%), B 1/60 (1.67\%), and C 1/58 observed trajectories
(1.72\%), with 57 observed failures and two unavailable C outcomes
(Table~\ref{tab:main}; Wilson intervals assume independent problem units, weakened
by related families). 

\subsection{Matched effects and uncertainty}
On the 58 complete pairs C$-$B is 0.00 points. The conventional exact two-sided
McNemar value is $p=1.0$, which carries no evidential value for equivalence.
Problem-level and whole-family bootstrap intervals collapse to $[0.00,0.00]$
because the observed paired vectors are identical; this is a degenerate artifact of
resampling the data, not precision about new theorems and not evidence of a zero
population effect. Table~\ref{tab:paired} also gives a separately labelled
Newcombe paired-score interval (Wilson marginals with an adjusted phi
coefficient~\cite{newcombe1998paired}), which conveys that materially different
population effects remain compatible with such sparse data; it assumes independent
pairs and is neither a cluster interval nor an equivalence bound. All estimates use
the complete cases appropriate to each contrast (C$-$A: 58 pairs; B$-$A: 60). Exact
and Holm-adjusted $p$ values are 1; all three planned comparisons stay in Holm
accounting. The bootstrap uses 10{,}000 resamples, seed 20271021; the complete C
contrasts cover 52 family clusters (54 in the full cohort). Resampling cannot repair selection bias from missing C outcomes. The three quantities answer different questions: the exact test has no discordant pairs to weigh; the bootstrap describes variability of the observed data under resampling, which is nil when every resampled vector coincides; and the score interval is a sensitivity analysis under an independence assumption, the only one of the three that indicates which population differences remain compatible with the data. A post-hoc, assumption-based sensitivity calculation for 58 independent pairs shows how little this design could detect: if the discordance probability were 0.2, 80\% exact power would require a true C$-$B difference of about 16.45 points, and if it were 0.05 or 0.1 no difference reaches 80\% exact power. This describes design sensitivity under stated assumptions; it is neither observed power nor family-adjusted power.

\begin{table}[tbp]
\caption{Paired success differences (percentage points). Exact two-sided McNemar
$p=1.0$ for all three contrasts; only C$-$B has zero discordant pairs (C$-$A and
B$-$A each have one discordant success). The cluster bootstrap (whole-family) is
primary; for C$-$B it is degenerate because the observed paired vectors are
identical. Score intervals are paired-score sensitivity intervals that assume
independent pairs; neither interval type is an equivalence bound.}
\label{tab:paired}
\centering\small
\begin{tabular}{lcrccc}
\toprule
Contrast & Pairs & Diff. & Cluster 95\% & Score 95\% & Exact $p$\\
\midrule
C$-$B & 58 & $+0.00$ & [$+0.00$, $+0.00$] & [$-6.83$, $+6.83$] & 1.0\\
C$-$A & 58 & $+1.72$ & [$+0.00$, $+5.56$] & [$-4.65$, $+9.14$] & 1.0\\
B$-$A & 60 & $+1.67$ & [$+0.00$, $+5.45$] & [$-4.50$, $+8.86$] & 1.0\\
\bottomrule
\end{tabular}
\end{table}

\subsection{Unavailable outcomes and sensitivity}
The frozen protocol planned both complete-pair inference and a full-denominator
missing-as-failure analysis, which treats C as 1/60 (1.67\%; Wilson [0.29,8.86]\%)
without altering its observed 1/58 rate. Both C-unavailable targets failed under
B. Neither unavailable trajectory is an observed proof failure, and neither was imputed or regenerated. Three quantities should not be conflated: the observed primary rate (1/58), the
missing-as-failure sensitivity rate (1/60), and the full-cohort identification
bound. Across the full cohort, C$-$B is zero if neither unobserved trajectory would
succeed and $+3.33$ points if both would; this $[0,+3.33]$ range is an
identification bound, not a sampling confidence interval, and is distinct from the
$\pm6.83$ independent-pair score interval and from the zero-width bootstrap interval.

The two absences have distinct documented mechanisms. The C1 prompt for
\texttt{lib\_a\_trig\_unit} has 9,653 input tokens, exceeding the
8,192-token cap before reserving 2,000 output tokens; the preflight check produced
no response. For \texttt{lib\_i\_bernoulli} the C1 prompt has 2,205 input tokens
but only a stream header survives, with zero durably observed generated tokens;
the cause is unknown and is not shown to be a context-limit failure. The no-replay
rule forbids regenerating either call with changed prompts, caps or seeds. Such
limits matter precisely because feedback lengthens inputs while verified successes
remain scarce.

\begin{figure}[tbp]
\centering
\includegraphics[width=.62\linewidth]{fig2.png}
\caption{Strictly verified theorem targets by experimental condition.
A verifies 0/60, B verifies 1/60, and C verifies 1/58 observed
trajectories (two unavailable). Both repair conditions solve the same
single target. These final totals do not establish a guided-repair advantage.}
\label{fig:verified-success}
\end{figure}

\subsection{Attempts, tokens and latency}
The ledger has 412 logical generation intents but 410 completed responses: 60
shared initials, 178 B repairs and 172 C repairs. C consumes 174 repair intents,
including one preflight rejection and one interrupted stream; an intent is not a
completed inference. Mean completed repairs per theorem are 2.97 (B) and 2.87 (C;
2.90 intents); both medians are three. Counting shared initials once, physical
completed calls consume 362,035 input and 130,401 output tokens and 11{,}233.6 s of
native generation; per-condition rows each include their initial and must not be
summed. C's diagnostic-bearing prompts are longer (Table~\ref{tab:res}). Call wall
time includes model loading, generation time does not; B always preceded C, cache
and hardware state were not randomized, and an unrelated MLX workload was observed,
so these are descriptive costs, not equal-cost treatment comparisons or causal
latency effects. The recovery elapsed record additionally includes an inactive
interval and is not continuous compute. More tokens and changed candidates
are not substitutes for more verified proofs.

\begin{table}[tbp]
\caption{Completed-call resource distributions (mean / median; loading included in
call wall time) and per-theorem means on the 58 complete pairs (230 completed
calls per condition, including the shared initial). Calls are not independent
theorem observations; missing costs are not imputed.}
\label{tab:res}
\centering\footnotesize\setlength{\tabcolsep}{3pt}
\resizebox{\linewidth}{!}{\begin{tabular}{lcccccc}
\toprule
Cond. & Calls & Input tok. & Output tok. & Call wall s & Gen.\ s & Per-theorem in/out tok.\\
\midrule
A & 60  & 215.5 / 215.0   & 464.4 / 198.0 & 31.8 / 16.3 & 27.2 / 11.8 & ---\\
B & 238 & 512.7 / 369.0   & 350.8 / 172.5 & 34.1 / 16.9 & 26.1 / 12.3 & 1980.9 / 1336.2\\
C & 232 & 1090.2 / 968.0  & 322.3 / 170.5 & 42.5 / 21.0 & 28.7 / 16.6 & 4353.5 / 1249.2\\
\bottomrule
\end{tabular}}
\end{table}

On the same 58 complete B/C theorem pairs, per-theorem wall time is
131.7 s mean / 68.6 s median for B versus 167.9 s / 84.8 s for C;
mean native generation time is 99.6 s versus 112.5 s. Matched mean
input/output token counts are 1980.9/1336.2 (B) versus 4353.5/1249.2
(C). The input-token ratio is approximately 2.2 (4353.5/1980.9) on these matched per-theorem totals, which include each condition's shared initial generation; it is a descriptive cost ratio, not an equalized budget or a causal efficiency effect. These observations describe unequal realized costs, not a
randomized or token-matched speed comparison.

%==================================================================
\section{Compiler-Error Taxonomy}
The frozen taxonomy assigns each failed initial candidate to one mutually
exclusive category in precedence order: parser/syntax, unknown name,
type/elaboration, tactic/goal, other. Of the 60 initial failures, 26 are
parser/syntax, 20 unknown-name, nine tactic/goal, four type/elaboration and one
other; this describes the model's starting errors, not mathematical difficulty,
and implies no ranking of categories by difficulty.
The only repaired success is in the unknown-name group (B and C each 1/20); all
other groups have zero verified repairs. Because C is unavailable in one parser and
one tactic/goal trajectory, its observed denominators there are 25 and eight.
Table~\ref{tab:errors} is descriptive; no
error-type superiority or subgroup-significance claim is made.

\begin{table}[tbp]
\caption{Primary initial-error taxonomy and observed repair success. Categories are
pre-specified and mutually exclusive; denominators are observed B/C outcomes.
Groups below five cases are descriptive.}
\label{tab:errors}
\centering\small
\begin{tabular}{lcccc}
\toprule
Initial error & Cases & B solved/obs. & C solved/obs. & C missing\\
\midrule
Other            & 1  & 0/1  & 0/1  & 0\\
Parser/syntax    & 26 & 0/26 & 0/25 & 1\\
Tactic/goals     & 9  & 0/9  & 0/8  & 1\\
Type/elaboration & 4  & 0/4  & 0/4  & 0\\
Unknown name     & 20 & 1/20 & 1/20 & 0\\
\bottomrule
\end{tabular}
\end{table}

Diagnostics illustrate practical obstacles (unavailable symbols, malformed syntax,
ill-typed applications, unsolved goals), but a message can locate a failed step
without supplying the correct lemma or argument. Unknown identifiers and unknown
theorem names are not reliably separable from wording alone, and tactic and
unsolved-goal messages can co-occur; an explicitly post-hoc message-tag artifact
keeps such overlapping labels without replacing the primary groups. No Lean
timeout was observed among completed candidates; the context-cap rejection is a
generation resource failure, not a compiler failure.

%==================================================================
\section{Exploratory Process Analysis: Candidate Repetition Under Diagnostics}
\emph{Post-hoc exploratory analysis; not a second primary endpoint.} After inspecting
the primary outcomes we asked whether diagnostics were associated with changes in
repeated proof strings. The outcome is byte-identical repetition of the preceding
candidate, not validity or mathematical progress. Attempts within a theorem are
dependent and B/C candidate histories diverge after the first repair, so the 178 and
172 repairs are not independent theorem units. These analyses generate hypotheses
and do not establish diagnostic comprehension or causal effects.
Table~\ref{tab:rep} summarizes all repetition counts.

\begin{table}[tbp]
\caption{Byte-identical repetition of the preceding candidate (exploratory).
Held-out test: B vs.\ C by repair stage. Development (15 problems, one seed): B, C
and the mismatched-feedback control C\_placebo. Unknown-name disappearance is shown
for the analysis in the text (B's diagnostics were visible only to the analyst).}
\label{tab:rep}
\centering\footnotesize\setlength{\tabcolsep}{3pt}
\resizebox{\linewidth}{!}{\begin{tabular}{llccc}
\toprule
Set & Measure & B & C & C\_placebo\\
\midrule
Test, all repairs      & repeats      & 90/178 (50.6\%) & 70/172 (40.7\%) & ---\\
Test, first repair (matched) & repeats & 17/58 (29.3\%) & 10/58 (17.2\%) & ---\\
Test, repairs 2--3     & repeats      & 73/118 (61.9\%) & 60/114 (52.6\%) & ---\\
Test                   & unknown name gone & 12/76 (15.8\%) & 14/70 (20.0\%) & ---\\
Dev, all repairs       & repeats      & 22/45 (48.9\%)  & 23/45 (51.1\%)  & 32/45 (71.1\%)\\
Dev                    & unknown name gone & 1/18 (5.6\%) & 3/19 (15.8\%) & 2/17 (11.8\%)\\
\bottomrule
\end{tabular}}
\end{table}

\begin{figure}[tbp]
\centering
\includegraphics[width=.58\linewidth]{fig3.png}
\caption{Post-hoc matched first-repair byte-identical repetition on
58 complete B/C pairs. Eight pairs repeat in both conditions, nine
only under B, two only under C, and 39 under neither. Exact
two-sided McNemar $p=0.0654$; the analysis is exploratory and
does not measure proof correctness.}
\label{fig:matched-repetition}
\end{figure}

\paragraph{All completed repairs.} B repeats its previous candidate byte-for-byte in
90/178 completed repairs (50.6\%), C in 70/172 (40.7\%). The repair-weighted
C$-$B difference is $-9.86$ points (95\% whole-family bootstrap CI
$[-19.77,-0.50]$; 10{,}000 resamples, seed 20271021). Although the interval barely
excludes zero, the contrast was chosen after outcome inspection, is one of several
exploratory analyses, and is not multiplicity-adjusted; we call it
\emph{suggestive}, not confirmatory evidence that feedback helps.

\paragraph{Matched first repair.} The first repair allows a directly matched
comparison because B and C share exact predecessor bytes. Among the same 58 complete
first-repair pairs, B repeats 17/58 (29.3\%) and C 10/58 (17.2\%): C$-$B is $-12.07$
percentage points (nine B-only repeats, two C-only, eight both, 39 neither).
The exploratory two-sided exact McNemar $p=0.0654$, whereas an independently
recomputed 10,000-resample, 52-family percentile bootstrap (seed 20271021)
gives $[-23.19,-1.69]$ points. The two procedures point in different directions
relative to conventional thresholds, and we report both rather than
choose the favorable one: this metric was chosen post hoc, the bootstrap is
sensitive to sparse clustered data, and no multiplicity correction was
performed. The earlier unequal-denominator $[-22.00,-0.86]$ interval is
superseded and not used.

\paragraph{Later repairs.} For repairs 2--3, B repeats 73/118 (61.9\%) and C 60/114
(52.6\%) on branch-specific predecessors; because the histories have diverged, these
cannot be read as matched causal comparisons, and the pooled contrast mixes
different histories. Crucially, fewer repeats produced no additional verified
proofs.

\begin{figure}[tbp]
\centering
\includegraphics[width=.62\linewidth]{fig4.png}
\caption{Exploratory byte-identical repetition by compiler-diagnostic
length quartile among 172 completed C repairs. Counts are
24/44, 16/42, 13/43 and 17/43 from shortest to longest.
Repeated attempts within the same theorem are dependent;
no causal diagnostic-length effect is inferred.}
\label{fig:diagnostic-quartiles}
\end{figure}

\paragraph{Diagnostic-length sensitivity (post hoc).}
From recorded C repair metrics, quartiles of diagnostic-token length contain
44, 42, 43, and 43 completed attempts, with byte-identical repeats in
24/44 (54.5\%), 16/42 (38.1\%), 13/43 (30.2\%), and 17/43
(39.5\%), respectively. Repetition is not monotone in message length;
diagnostic length can covary with task difficulty, candidate history,
and failure type. These descriptive quartiles neither identify diagnostic
comprehension nor change the verified-success result.

\paragraph{Name removal is not evidence of diagnostic use.} An immediately
preceding unknown name disappears in 12/76 (15.8\%) eligible B repairs and 14/70
(20.0\%) C repairs. B never received the diagnostic, so blind name removal is a
chance-edit baseline; the small C excess does not show diagnostic use. Deleting a
symbol, shortening a message list or changing primary-error category does not
certify a valid proof step.

\paragraph{Edit distance and error-stage transitions (archived post hoc).} Two archived supplementary analyses of the same completed repairs, run without new inference, tell the same story. Mean normalized character edit distance between consecutive candidates is 0.1109 for B ($n=178$) and 0.1176 for C ($n=172$), with a cluster interval for C$-$B of $[-0.0207,+0.0327]$; the frozen-tokenizer analogue is 0.1172 versus 0.1245, interval $[-0.0236,+0.0360]$. Classifying each repair by the declared stage order of its primary error category (parser, unknown name, type, tactic), B's repairs advanced to a later stage in 11/178 (6.2\%), stayed in the same category in 161 (90.4\%), regressed in 4 (2.2\%), verified in 1 and were unranked in 1; C's advanced in 14/172 (8.1\%), stayed in 146 (84.9\%), regressed in 7 (4.1\%), verified in 1 and were unranked in 4. Fewer messages or a stage change do not certify a correct proof step, and no test is applied.

\paragraph{Development placebo.} The later post-hoc control is separate from the
three-case pre-test pilot. It uses fresh shared initials on 15 development problems
in five families, one seed, and real diagnostics donated from different
problems/families (only filename metadata normalized). All four arms (A, B, C,
C\_placebo) verify 0/15. Over 45 completed repairs per repair arm, repetition is
22/45 (48.9\%) for B, 23/45 (51.1\%) for C and 32/45 (71.1\%) for C\_placebo.
Real C therefore does \emph{not} show lower repetition than B on this set (it is
slightly higher), unlike the held-out aggregate. The higher placebo repetition is
compatible with several explanations, including confusion from irrelevant text or
copying under uninformative prompts; it does not isolate diagnostic comprehension.
The unknown-name differences (C 3/19 vs.\ placebo 2/17) rest on very small
denominators and support no inference. The placebo neither validates nor refutes a
diagnostic-content mechanism.

\paragraph{Limitations.} The placebo has 15 problems, five families, one seed and no hypothesis test. Branch order, random streams, unequal feedback lengths, repeated observations, donor-message choice and post-hoc selection prevent causal attribution; no replay or imputation was run, and both unavailable C outcomes stay unavailable.

%==================================================================
\section{Qualitative Repair Cases}
Five authentic cases, selected post hoc for outcome and failure diversity,
illustrate that a modified proof is not necessarily a repaired proof. These
are not five successes: the experiment verifies one distinct held-out target.
Compiler messages associated with B are displayed \emph{for analysis only};
blind repair never received them. The descriptions below summarize observed
changed lines and verification records, not stand-alone compilable Lean code.
The full original candidate bytes, hashes and diagnostics remain in the ledger.

\begin{table}[tbp]
\caption{Authentic final repair cases (selected post hoc). ``Identical'' means
byte-identical to the immediate predecessor; diagnostic exposure applies to C only.}
\label{tab:cases}
\centering\footnotesize\setlength{\tabcolsep}{4pt}
\begin{tabular}{p{2.3cm}p{1.0cm}p{3.2cm}p{4.0cm}}
\toprule
Target & Final & Preceding failure & Observed change and outcome\\
\midrule
\texttt{lib\_a\_exp\_add} & B1 & Unknown \texttt{Real.ofReal} & Removed invalid intermediate \texttt{have}; retained \texttt{rw [Real.exp\_add]}; \textbf{verified}.\\
\texttt{div\_007} & C3 & Unknown \texttt{even\_add} & Identical candidate; not verified.\\
\texttt{lib\_a\_trig\_}\allowbreak\texttt{unit} & B3 & Unsolved trigonometric equality goal & Changed parentheses/expression in an intermediate statement; not verified.\\
\texttt{div\_008} & C3 & Failed \texttt{split} tactic & Inserted \texttt{by rw [mul\_one,} \texttt{hk, mul\_one]} elsewhere; retained failing tactic; not verified.\\
\texttt{div\_011} & C3 & Integer term supplied where a divisibility proof was required & Identical candidate; not verified.\\
\bottomrule
\end{tabular}
\end{table}

In \texttt{lib\_a\_exp\_add}, blind repair removed a failed intermediate
use of \texttt{Real.ofReal} and proved the identity using the library lemma
\texttt{Real.exp\_add} without seeing the diagnostic; C also verified this same
theorem. This single genuine success shows that a compiler diagnostic is not always
necessary for a repair to succeed, and it is the only evidence of verified repair
in the study. In contrast, \texttt{div\_007} and \texttt{div\_011} show exact repeats
despite authentic feedback, while \texttt{lib\_a\_trig\_unit} (a completed B3
attempt; the unavailable outcome for this target is C1) and
\texttt{div\_008} change candidate bytes but remain rejected. Together, the two guided
repeats and two changed-but-failing candidates show the limits of string-change
evidence: neither an identical nor a modified candidate says whether the proof
moved closer to acceptance. The actual
revision to \texttt{div\_008} was outside the original cropped five-line display;
no fabricated before/after proof is supplied here. Case selection is
illustrative, not an estimate of failure-mode prevalence.

%==================================================================
\section{Discussion, Implications and Threats to Validity}
\paragraph{Three evidence levels.} The results look different depending on what is
measured. (i)~At the level of \emph{candidate strings}, offering diagnostics was
associated with fewer byte-identical retries on the held-out set (exploratory,
unadjusted, history-dependent), although not on the development set. (ii)~At the
level of \emph{compiler-error descriptions and transitions}, the primary error category stayed the same in most repairs under both arms (90.4\% of B and 84.9\% of C repairs), later-stage advancement was uncommon (6.2\% vs.\ 8.1\%), and unknown names also disappeared under blind repair (15.8\% vs.\ 20.0\%), so such transitions do not certify use of the diagnostic. (iii)~At the level of
\emph{original-theorem kernel verification}, the frozen primary endpoint, B and C
verify the same single theorem. Only the third level answers RQ1, and it shows no
observed advantage. The methodological point is that a study reporting only the
first two levels could have reported an apparent benefit. The result is bounded:
under this model, snapshot, seed and three-repair cap, diagnostics did not produce
additional verified targets.

\paragraph{Competing explanations.} The extremely low observed verification rate leaves little empirical information for estimating between-condition differences. This sparse-success regime is compatible with several explanations, including limited model capability, task difficulty and ineffective use of feedback, but the study does not establish which mechanism predominates. Further candidates include a 7B 4-bit model lacking the relevant library knowledge or planning ability; longer guided prompts that may dilute the task or exceed the context cap; a blind-feedback text that may not be neutral; B-before-C ordering and condition-specific random streams; and a repetition metric that reflects surface behaviour unrelated to proof progress. A sparse-success regime may limit sensitivity, but we do not treat this as an established cause. None is tested here.

\paragraph{Feasibility and cost.} Guided prompts raised matched per-theorem input tokens from 1980.9 to 4353.5 and made one trajectory impossible under the 8,192-token cap; for bounded-context local models, diagnostic length is a feasibility constraint, not only a cost.

\paragraph{Construct validity.} The checker guards the fixed type: a candidate cannot
weaken the conclusion, add a premise, change a quantifier or use unauthorized
axioms and still count. Kernel acceptance does not show that the type captures the
English. Historical coding-AI assessments judged 94 of 96 original references
faithful, one not faithful (\texttt{mod\_001}, missing a positive-modulus premise)
and one uncertain (\texttt{ind\_001}, ambiguous endpoints); both were corrected but
excluded for provenance. The retained 75/75 AI eligibility decisions are not an
independently validated accuracy claim (Appendix~A); autonomous
natural-language-to-type generation was never measured. Repetition is a surface
measure and not mathematical progress.

\paragraph{Internal and statistical validity.} B and C start identically with equal
repair opportunities, but prompts differ in length and condition-specific random
stream; B precedes C; costs are not controlled for tokens or machine state; two C
outcomes are unavailable; and the protocol was frozen internally, with process
outcomes analysed post hoc. With one solved theorem the paired vectors have zero
discordance and bootstrap intervals collapse. Family clustering, sparse subgroup
counts, unequal repair denominators, one seed and several unadjusted exploratory
comparisons further constrain inference; the wider paired-score interval and the
missing-outcome bounds are reported rather than hidden.

\paragraph{External and semantic validity.} The authoring/adjudicating assistant may
share biases with dataset construction; separate Phi reviewers proved unreliable and
bounded finite probes cannot exhaust informal meaning. The data include
formal-library adaptations and are not contamination-free; curated families are human
groupings, not independence guarantees. One 4-bit model, one host and one
Lean/Mathlib snapshot were used, and the model may predate library interfaces.
Results should not be generalized to trained repair systems, larger models, other
domains or unconstrained formalization.

\paragraph{Prospective implications.} The study motivates predeclaring verified
success as the decisive endpoint, reporting string changes and error transitions as
secondary process measures, retaining paired initial states, measuring missingness
rather than silently rerunning, and costing diagnostic prompt length. A natural
follow-up, outside the frozen study and not performed here, would use a model or
prompt with measurable baseline success, multiple seeds, a placebo arm on held-out
data, token-matched blind prompts (e.g.\ padded with irrelevant text), and a
concise-diagnostic variant.

%==================================================================
\section{Reproducibility and Authorization Boundary}
The internal freeze binds protocol, dataset and model inventory to recorded hashes
(protocol SHA-256 \path{a032110d7afbc367c2daa820bd17511c6ff5dba2d1164e0130ea5d1c4b693317};
dataset freeze \path{de0850197b821ca4c519737c4cffb16d78e2a329d15be6c7f3c6a7f5487db67f};
model inventory \path{24d68abbe5e1434e5263ad123cca517c2c1799c7465e104e83f6cb02baeac4f9}),
and the ledger preserves raw generations, compiler JSON, verification decisions,
tokens, timings and checkpoints. The pre-held-out source commit, launch record, evidence hashes and release
build manifest bind execution to the frozen inputs. Reconstruction re-checks
every completed token journal against its raw response, re-tokenizes prompts
offline, and deterministically regenerates metrics; both successful
candidates were reverified against exact targets and allowed axioms. Raw
evidence whitespace is preserved, with editable-source whitespace checked
separately.

Two reproducibility levels should be distinguished. \emph{Public aggregate
reanalysis} is supported by the released redacted views. \emph{Exact-byte
reproduction} is not: released evidence redacts private paths, while recorded prompt
hashes, derived seeds and tokenization refer to the private original bytes, so exact
prompt replay and validation require the preserved originals. Complete public
byte-for-byte replay is not established. The freeze is internal, not public
preregistration, and reference screening was AI-assisted and is not described as
human review. Repository records report a reconciliation of the post-hoc analyses against the original private records (38 checks, zero mismatches); we did not re-execute it. Appendix~B states the authorization and release limits.

%==================================================================
\section{Conclusion}
In this fixed-theorem, local-model pilot, genuine compiler diagnostics did not
verify any theorem that blind retries did not: B and C verify the same single
target (1/58 on complete pairs). Such sparse outcomes can support no claim of
equivalence, and the two unavailable C outcomes, the wide paired-score interval and
the single-seed setting narrow the claim further. Because the blind comparator does not isolate diagnostic content from prompt length, random streams or execution order, this is a comparison of two repair procedures, not a test of a mechanism.

Post hoc, guided retries repeated exact candidate bytes less often than blind
retries on the held-out set (70/172 vs.\ 90/178), but this exploratory difference
did not translate into successful proofs, the development placebo is ambiguous, and
blind retries also remove unknown names. The evaluation lesson is to
report verified progress, superficial candidate change and diagnostic-prompt cost as
distinct outcomes. The scope is a fixed-type proof-synthesis stage with one model
and seed; the study neither evaluates autonomous English-to-Lean translation nor
establishes a general limitation of compiler feedback.

%==================================================================
\appendix
\section{Dataset Provenance and Semantic Audit}
From 96 prepared candidates we exclude one exact duplicate and 20 nonduplicate
records lacking adequate item-level provenance, leaving 75 records: 16
textbook-derived statements~\cite{levin2021dmoi} (CC BY-SA 4.0 attribution) and 59
adaptations of pinned formal-library statements (Apache-2.0 notices). Citations,
original source bytes, license notices and item-level attribution are retained;
these are evidence-based AI-assisted source assessments, not publisher permissions
or independent legal clearance. Two original records were corrected (a restored
positive-modulus premise; explicit summation endpoints), reverified and re-reviewed
with Phi-4, and remain excluded for provenance; corrections were not used to inflate
eligible counts.

The replacement policy uses explicit AI adjudication of quantifiers, domains,
assumptions, strength, endpoints and equivalent encodings. Phi-3.5 and Phi-4 audit
responses are retained, including schema failures and unsupported objections; the
authoring coding assistant adjudicated the exact selected pairs. This is not
author-independent human review, and unreliable raw AI output is not counted as a
faithful vote. The reviewer models were demonstrably fallible (one claimed
$1(1+1)$ was odd). All 96 original reference types have strict Lean evidence, and
bounded probes cover 26,467 accepted finite assignments (88 exact predicates, eight
approximate numerical probes); these supplement, not establish, English equivalence.
Whole binomial, linear-solvability and modular-power families were grouped before
primary-model test access. No held-out outcome was used to select prompts or
settings, and the freeze is internal, not external preregistration.

\begin{table}[tbp]
\caption{Frozen composition. Category, difficulty and provenance rows are separate
partitions, not additive totals. Development/test occupy 5/54 curated families with
no family overlap.}
\label{tab:comp}
\centering\small
\begin{tabular}{lrrr}
\toprule
Partition & Dev. & Test & Total\\
\midrule
Category: arithmetic identity & 2 & 11 & 13\\
Category: divisibility        & 3 & 17 & 20\\
Category: induction           & 6 & 8  & 14\\
Category: inequality          & 1 & 15 & 16\\
Category: modular arithmetic  & 3 & 9  & 12\\
Difficulty: 1                    & 2 & 2  & 4\\
Difficulty: 2                    & 10 & 57 & 67\\
Difficulty: 3                    & 2 & 1  & 3\\
Difficulty: 4                    & 1 & 0  & 1\\
Textbook / CC BY-SA 4.0       & 8 & 8  & 16\\
Formal library / Apache-2.0   & 7 & 52 & 59\\
\midrule
Total unique statements       & 15 & 60 & 75\\
\bottomrule
\end{tabular}
\end{table}

\begin{table}[tbp]
\caption{AI-assisted reference faithfulness; not autonomous translation accuracy or
human-gold accuracy. Original 96-item labels precede correction, source exclusion
and final adjudication.}
\label{tab:faith}
\centering\small
\begin{tabular}{lccc}
\toprule
Evidence population & Faithful & Not faithful & Uncertain/invalid\\
\midrule
Coding-AI original references (96)   & 94 & 1  & 1 uncertain\\
Phi-4 original responses (96)        & 72 & 11 & 13 invalid\\
Selected AI-eligible references (75) & 75 & 0  & 0 uncertain\\
Independent human labels             & 0  & 0  & Not measured\\
\bottomrule
\end{tabular}
\end{table}

\section{Authorization and Release Limits}
The author reports that supervisor approval of AI-assisted semantic review and free
local MLX substitution was given on 28 October 2026; the decision document and
signature are not included here, and the approval is not independently verified.
Implementation choices were frozen under the author's execution authorization. No
publisher permission or human review was fabricated. AI tools assisted construction,
adjudication and inference; they are not credited as authors. Public evidence is a
privacy-redacted view: paths in diagnostics, prompts and metadata are replaced in
exported copies, while recorded hashes, token counts and seeds still describe the
original private bytes, so the original frozen validators must not be applied to
redacted views. Original complete test logs and historical exploratory evidence
remain local. This is a disclosed public reproducibility limit, not a revision of
the experiment.

%==================================================================
\begin{thebibliography}{20}
\bibitem{azerbayev2023proofnet} Azerbayev, Z., Piotrowski, B., Schoelkopf, H., Ayers, E.W., Radev, D., Avigad, J.: ProofNet: Autoformalizing and formally proving undergraduate-level mathematics (2023), \url{https://arxiv.org/abs/2302.12433}
\bibitem{chen2024selfdebug} Chen, X., Lin, M., Sch{\"a}rli, N., Zhou, D.: Teaching large language models to self-debug. In: ICLR (2024), \url{https://arxiv.org/abs/2304.05128}
\bibitem{demoura2021lean4} de Moura, L., Ullrich, S.: The Lean 4 theorem prover and programming language. In: Automated Deduction -- CADE 28. LNCS, vol. 12699, pp. 625--635. Springer (2021). \doi{10.1007/978-3-030-79876-5_37}
\bibitem{first2023baldur} First, E., Rabe, M.N., Ringer, T., Brun, Y.: Baldur: Whole-proof generation and repair with large language models. In: ESEC/FSE (2023), \url{https://arxiv.org/abs/2303.04910}
\bibitem{huang2024selfcorrect} Huang, J., Chen, X., Mishra, S., Zheng, H.S., Yu, A.W., Song, X., Zhou, D.: Large language models cannot self-correct reasoning yet. In: ICLR (2024), \url{https://arxiv.org/abs/2310.01798}
\bibitem{jiang2023dsp} Jiang, A.Q., et al.: Draft, sketch, and prove: Guiding formal theorem provers with informal proofs. In: ICLR (2023), \url{https://arxiv.org/abs/2210.12283}
\bibitem{levin2021dmoi} Levin, O.: Discrete mathematics: An open introduction (2021), \url{https://discrete.openmathbooks.org/dmoi3.html}, third edition, fifth printing. Source states CC BY-SA 4.0; source and license evidence archived
\bibitem{madaan2023selfrefine} Madaan, A., et al.: Self-Refine: Iterative refinement with self-feedback. In: NeurIPS (2023), \url{https://arxiv.org/abs/2303.17651}
\bibitem{mathlib2020} The mathlib Community: The Lean mathematical library. In: Proc. 9th ACM SIGPLAN Int. Conf. on Certified Programs and Proofs, pp. 367--381. ACM (2020). \doi{10.1145/3372885.3373824}
\bibitem{mlxlm} MLX LM: Large language model inference and fine-tuning for Apple silicon. \url{https://github.com/ml-explore/mlx-lm}, installed version 0.32.2; accessed 27 October 2026
\bibitem{newcombe1998paired} Newcombe, R.G.: Improved confidence intervals for the difference between binomial proportions based on paired data. Statistics in Medicine 17(22), 2635--2650 (1998). \doi{10.1002/(SICI)1097-0258(19981130)17:22<2635::AID-SIM954>3.0.CO;2-C}
\bibitem{qwenmlx} MLX conversion of Qwen2.5-7B-Instruct, four-bit. \url{https://huggingface.co/mlx-community/Qwen2.5-7B-Instruct-4bit}, snapshot c26a38f6a37d0a51b4e9a1eb3026530fa35d9fed; model card and Apache-2.0 license evidence archived 27 October 2026
\bibitem{qwen2024report} Qwen Team (Yang, A., et al.): Qwen2.5 technical report (2024), \url{https://arxiv.org/abs/2412.15115}
\bibitem{shinn2023reflexion} Shinn, N., Cassano, F., Berman, E., Gopinath, A., Narasimhan, K., Yao, S.: Reflexion: Language agents with verbal reinforcement learning. In: NeurIPS (2023), \url{https://arxiv.org/abs/2303.11366}
\bibitem{wang2026april} Wang, E., Chess, S., Lee, D., Ge, S., Mallavarapu, A., Ilin, V.: Learning to repair Lean proofs from compiler feedback. arXiv preprint (2026), \url{https://arxiv.org/abs/2602.02990}
\bibitem{xin2024deepseekprover} Xin, H., et al.: DeepSeek-Prover: Advancing theorem proving in LLMs through large-scale synthetic data (2024), \url{https://arxiv.org/abs/2405.14333}
\bibitem{yang2023leandojo} Yang, K., Swope, A., Gu, A., Chalamala, R., Song, P., Yu, S., Godil, S., Prenger, R.J., Anandkumar, A.: LeanDojo: Theorem proving with retrieval-augmented language models. In: Advances in Neural Information Processing Systems, vol.~36 (2023), \url{https://arxiv.org/abs/2306.15626}
\bibitem{zhang2026beyond} Zhang, K., Gallardo Candela, P., Murthy, S., Xie, Y., Wang, Z., Raissi, M.: Beyond compilation: Evaluating faithful natural-language-to-Lean statement formalization (2026), \url{https://arxiv.org/abs/2606.31002}
\bibitem{zheng2022minif2f} Zheng, K., Han, J.M., Polu, S.: miniF2F: A cross-system benchmark for formal olympiad-level mathematics. In: ICLR (2022), \url{https://arxiv.org/abs/2109.00110}
\bibitem{li2026hindsight} Li, J., Li, S.:
Repair, Retry, and the Cost of Hindsight: Verifier Feedback in Formal Theorem
Proving. arXiv:2610.22698 (2026),
\url{https://arxiv.org/abs/2610.22698}
\end{thebibliography}

\end{document}